\documentclass[a4paper,12pt]{article}

\usepackage[utf8]{inputenc}
\usepackage[portuguese]{babel}
\usepackage{fancyhdr}
\usepackage{lastpage}
\usepackage{amsmath}
\usepackage{amsthm}
\usepackage{amssymb}
\usepackage[portuguese,algoruled,lined,linesnumbered]{algorithm2e}
\usepackage{geometry}
\usepackage{hyperref}
\geometry{a4paper, margin=1in}


% ---- PREENCHA ----
\newcommand{\Lista}{7}
\newcommand{\Nome}{Leonardo Heidi Almeida Murakami}
\newcommand{\NUSP}{11260186}
% -----------------------------

\pagestyle{fancy}
\fancyhf{}
\fancyhead[L]{\Nome}
\fancyhead[R]{NUSP \NUSP}
\fancyfoot[L]{Lista \Lista}
\fancyfoot[C]{MAC0338}
\fancyfoot[R]{Página \thepage\ de \pageref{LastPage}}

\begin{document}

\section*{Exercício 3}

\subsection*{Análise}
A afirmação é falsa.

Sabemos que a propriedade fundamental da MST relacionada a ciclos estabelece que a aresta de custo máximo em qualquer ciclo não pode pertencer à MST (supondo custos distintos). A pergunta, porém, trata da aresta de custo mínimo, para a qual não há tal garantia.

Provaremos isso por contraexemplo, construindo um grafo onde a aresta de custo mínimo de um ciclo $C$ não pertence à MST.

\subsection*{Contraexemplo}
Considere o grafo $G=(V, E)$ com 4 vértices $V = \{A, B, C, D\}$ e 5 arestas:
\begin{itemize}
    \item $c(A, D) = 2$
    \item $c(D, B) = 3$
    \item $c(A, B) = 5$
    \item $c(B, C) = 6$
    \item $c(C, A) = 7$
\end{itemize}
Como o grafo é conexo e todos os custos são distintos, existe uma única MST.

Considere o ciclo $C = A-B-C-A$. As arestas deste ciclo são $(A, B)$ com custo 5, $(B, C)$ com custo 6 e $(C, A)$ com custo 7. Logo, a aresta de custo mínimo em $C$ é $(A, B)$.

Construindo a MST de $G$ pelo algoritmo de Kruskal:
\begin{enumerate}
    \item Adiciona $(A, D)$ (custo 2). Arestas da MST: $\{(A, D)\}$.
    \item Adiciona $(D, B)$ (custo 3). Arestas da MST: $\{(A, D), (D, B)\}$.
    \item Ao considerar $(A, B)$ (custo 5), observamos que ela forma o ciclo $A-D-B-A$. Como $(A, B)$ é a aresta de maior custo neste ciclo, ela é rejeitada.
    \item Adiciona $(B, C)$ (custo 6). Arestas da MST: $\{(A, D), (D, B), (B, C)\}$.
    \item Ao considerar $(C, A)$ (custo 7), observamos que ela forma o ciclo $A-D-B-C-A$. Como $(C, A)$ é a aresta de maior custo neste ciclo, ela é rejeitada.
\end{enumerate}
Portanto, a MST de $G$ é $T = \{(A, D), (D, B), (B, C)\}$, com custo total 11.

\subsection*{Conclusão}
A aresta $(A, B)$ é a aresta de custo mínimo no ciclo $C$, mas não pertence à MST de $G$. Isso ocorre porque, apesar de ser mínima em $C$, ela é máxima em outro ciclo ($A-D-B-A$), sendo portanto excluída pela propriedade de ciclo.

Concluímos que a afirmação é falsa.

\newpage

\section*{Exercício 5}

\subsection*{Prova}
A afirmação é verdadeira. Provaremos usando a Propriedade do Corte das MSTs.

Seja $G = (V, E)$ o grafo e $T$ uma MST que contém a aresta $e = (u, v)$.

\begin{enumerate}
    \item \textbf{Construção do Corte:}
    Ao remover a aresta $e$ da árvore $T$, obtemos exatamente dois componentes conexos. Seja $S$ o conjunto de vértices do componente que contém $u$ e $V-S$ o conjunto de vértices do componente que contém $v$. Esta partição $(S, V-S)$ define um corte no grafo $G$, e a aresta $e$ cruza este corte por construção.
    
    \item \textbf{Prova por Contradição:}
    Suponha, por absurdo, que $e$ não seja de peso mínimo entre as arestas que cruzam o corte $(S, V-S)$. Então existe uma aresta $e' = (x, y)$ em $G$ tal que $e'$ cruza o corte e $w(e') < w(e)$.

    \item \textbf{Argumento de Troca:}
    Se adicionarmos $e'$ a $T$, criamos um ciclo, pois $T$ já é uma árvore. Este ciclo necessariamente cruza o corte $(S, V-S)$ duas vezes: uma pela aresta $e'$ e outra pela aresta $e$, que é a única ponte entre $S$ e $V-S$ em $T$. Logo, o ciclo contém ambas $e$ e $e'$.
    
    \item \textbf{Construção de $T'$:}
    Defina $T' = (T - \{e\}) \cup \{e'\}$. Como removemos uma aresta do ciclo, $T'$ continua sendo uma árvore geradora. Comparando os pesos:
    $$ w(T') = w(T) - w(e) + w(e') < w(T) $$
    pois $w(e') < w(e)$. Mas isto contradiz o fato de $T$ ser uma MST.

    \item \textbf{Conclusão:}
    Logo, nossa hipótese de que existe $e'$ com peso menor que $e$ cruzando o corte é falsa. Portanto, $e$ tem peso mínimo entre todas as arestas que cruzam $(S, V-S)$.
\end{enumerate}

\end{document}